% TOP_PROBLEM: {"id": 3700021, "title": "Small components of subharmonic level sets", "category_id": 8, "category": {"id": 8, "name": "analysis", "display_name": "Analysis", "description": "Limits, continuity, calculus, and function theory.", "slug": "analysis", "order_index": 8, "created_at": "2026-07-31T15:26:25.671Z"}, "status": "partially_solved", "problem_number": "AMR-036-0021", "set_id": 13, "statusQualification": "Uses open horizontal sections inside the source's open square; discarding finitely many indices ensures the distinguished component exists."}
% FIRST_POSTED: 2026-10-01
% ENTRY_KIND: statement_only
% UnsolvedMath ID 3700021; source catalogue code AMR-036-0021
% !TeX program = lualatex
% Definition and sources only; this file is not an LLM solution attempt.
\documentclass[11pt,a4paper]{article}
\usepackage[margin=25mm]{geometry}
\usepackage{fontspec}
\setmainfont{lmroman10-regular.otf}[BoldFont=lmroman10-bold.otf,
 ItalicFont=lmroman10-italic.otf,BoldItalicFont=lmroman10-bolditalic.otf]
\usepackage{amsmath,amssymb,mathtools}
\usepackage{unicode-math}
\setmathfont{latinmodern-math.otf}
\AtBeginDocument{\let\setminus\smallsetminus}
\usepackage{xurl,hyperref}
\hypersetup{colorlinks=true,urlcolor=blue}
\setcounter{secnumdepth}{0}
\setlength{\parindent}{0pt}
\setlength{\parskip}{5pt}
\setlength{\emergencystretch}{3em}
\begin{document}
\section*{Small components of subharmonic level sets}\label{3700021}
{\small UnsolvedMath ID 3700021 · AMR-036-0021 · Analysis · AMR Open Problem Lists}
\subsection{Definitions and mathematical statement}
Let $S=\{x+iy:|x|<1,\ |y|<1\}$ be the open square.
A real-valued function on $S$ is subharmonic if it is upper
semicontinuous and its value at the center of every closed
disk contained in $S$ is at most its average over the
boundary circle.
Consider subharmonic functions $u_k$ satisfying uniform
convergence on the entire square:
\[
             \sup_{x+iy\in S}|u_k(x+iy)-x|\longrightarrow0.
\]
This is stronger than uniform convergence on compact subsets.

Put $D_k=\{z\in S:u_k(z)<0\}$. For sufficiently large
$k$, the point $-1/2$ belongs to this open set; let
$D_k^*$ be its connected component. Is it possible that,
after discarding finitely many terms, for every $k$ and
every $t\in(-1,1)$,
\[
       (D_k\setminus D_k^*)\cap
             \{x+it:-1<x<1\}\ne\varnothing?
\]
In words, the negative components other than the component
containing $-1/2$ must collectively meet every horizontal
cross-section of the square.

The intersecting component and the intersection point may
depend on the height $t$. There is no requirement that one
single small component cross every height. Uniform convergence
forces the parts outside the large negative component to
lie increasingly close to the vertical zero line of the
limit function; the question asks whether subharmonicity
nevertheless permits them at every height simultaneously.
\subsection{Short English statement}
Can extra negative components of subharmonic approximations to the real coordinate collectively reach every horizontal level?
\subsection{Sources}
\begin{itemize}
\item[S1] ULAM.AI and the UnsolvedMath Contributors, supplied problems.json, record 3700021 (AMR-036-0021), AMR Open Problem Lists. Catalogue identity and source transcription; CC BY 4.0.
\url{https://huggingface.co/datasets/ulamai/UnsolvedMath}.
\item[S2] Eremenko - My favorite unsolved problems, Potential theory, Problem 2. Original source indexed by AMR-036-0021.
\url{https://www.math.purdue.edu/~eremenko/uns1.html}.
\item[S3] A. Eremenko, Problems of potential theory arising in value distribution, Problem 2. Author-maintained primary problem note.
\url{https://www.math.purdue.edu/~eremenko/dvi/potvd.pdf}.
\end{itemize}
\end{document}
